% TikZ 图：向下运行的倾斜传送带初始阶段
% Day 07：连接体、斜面与摩擦综合
\begin{tikzpicture}[scale=0.98,>=Latex,line cap=round,line join=round]
  \def\ang{22}

  % 带面及其运动方向。
  \draw[thick] (0,0) -- (5.70,2.31);
  \draw[thick] (0,-0.28) -- (5.70,2.03);
  \foreach \x in {0.60,1.45,2.30,3.15,4.00}
    \draw[->,blue,thick] (\x,{0.405*\x-0.14}) -- ++(202:0.48);

  \draw[->,ultra thick,blue] (5.20,2.58) -- ++(202:1.15)
    node[midway,above=4pt] {$v_{\text{带}}$};

  % 物块下表面与带面贴合。
  \coordinate (P) at (2.36,0.96);
  \coordinate (C) at ($(P)+(112:0.36)$);
  \begin{scope}[shift={(P)},rotate=22]
    \draw[fill=white,thick] (-0.55,0) rectangle (0.55,0.72);
    \node at (0.27,0.48) {$m$};
  \end{scope}

  % 初始阶段的三个真实力，均从物块中心出发。
  \draw[->,very thick,red] (C) -- ++(202:1.12)
    node[below left=1pt] {$f_k$};
  \draw[->,very thick,blue] (C) -- ++(0,-1.28)
    node[below right=1pt] {$mg$};
  \draw[->,very thick,blue] (C) -- ++(112:1.02)
    node[above left=1pt] {$N$};

  % 倾角与阶段说明。
  \draw[thin] (0,0) -- (1.25,0);
  \draw (0.72,0) arc[start angle=0,end angle=22,radius=0.72];
  \node at (0.98,0.17) {$\theta$};
  \node[font=\footnotesize,align=center] at (3.15,-0.72)
    {初始时物块相对带面向上滑\\
     滑动摩擦力沿斜面向下};
\end{tikzpicture}
